\chapter{Serie}

\section{Serie Numeriche}

\defn{Serie Numerica}{
    Sia $a_{n}$ una successione di numeri reali, si chiama \textbf{serie numerica di termine generale $a_{n}$} la seguente successione detta delle \textbf{somme parziali}.
    \begin{align*}
        & S_{1}=a_{1}\\
        & S_{2} = S_{1}+a_{2}=a_{1}+a_{2} \\
        & S_{3} = S_{2}+a_{3}=a_{1}+a_{2} +a_{3} \\
        & \vdots\\
        & S_{n} = S_{n-1}+a_{n} = a_{1}+a_{2}+a_{3}+\dots+a_{n}
    \end{align*}
    \[ S_{n}=\sum_{k=1}^{n} a_{k} \]
}

\subsection{Convergenza o divergenza}

\defn{Serie convergente}{
    Diremo che la serie \textbf{converge alla somma $S\in \mathbb{R}$} se $S_{n}$ è convergente a
    \[ \lim_{ n \to \infty }S_{n}=S \]
    e $S$ si indica con 
    \[ S=\sum_{n=1}^{+\infty} a_n \]
}

\defn{Serie divergente}{
    Diremo che la serie \textbf{diverge a $\pm \infty$} se
    \[ \lim_{ n \to \infty } S_{n}=\pm \infty \]
    rispettivamente:
    \[ \sum_{n=1}^{+\infty} a_{n}=\pm \infty \]
}

\defn{Serie Regolare}{
    Se la serie diverge o converge, diremo che è una serie regolare.
}

\defn{Serie Non Regolare}{
    Se: 
    \[ \lim_{ n \to \infty } S_{n} = \not \exists \]
    diremo che non è regolare.
}

D'ora in poi una serie di termine generale $a_{n}$ si indicherà con:
\[ \sum_{n=1}^{+\infty} a_{n} \]
studiare il \textbf{carattere} delle serie significa studiare se la serie:
\begin{itemize}
    \item converge
    \item diverge
    \item oscilla
\end{itemize}

\subsection{Esempi di tipi di Serie}

\subsubsection{Serie Geometrica}

\defn{Serie Geometrica}{
    Sia $q \in \mathbb{R}$.
    La serie geometrica di ragione $q$:
    \[ \sum_{n=0}^{+\infty} q^n = 1 + q+ q^2 + \dots \]
    Si chiama così perché ricorda la progressione geometrica.
}

\ex{
    Progressione geometrica:
    \[
    S_{n}=\begin{dcases}
    \frac{1-q^{n+1}}{1-q} \qquad q\neq 1 \\
    \text{Sommo sempre 1 } \quad q=1
    \end{dcases}
    \]
    \[
    \lim_{ n \to \infty } S_{n}= \begin{dcases}
    \frac{1}{1-q} =S \quad &|q|<1 \\
    \not \exists & q \le -1 \\
    +\infty & q\ge 1
    \end{dcases}
    \]
}

\ex{
    \[
    \begin{drcases}
    & 0,\overline{9}=0,9+0,09+0,009+0,0009+\dots = \\
    & =\frac{9}{10} + \frac{9}{100}+\frac{9}{1000} + \dots + \frac{9}{10^n}= \\
    & = \frac{9}{10}\left( 1 +\frac{1}{10} +\frac{1}{100} + \dots \right) = \frac{9}{10}\left( \frac{1}{1-\frac{1}{10}} \right)=\frac{9}{10} \left(\frac{10}{9}\right) = 1
    \end{drcases}0,\overline{9}=1
    \]
}

\subsubsection{Serie Armonica}
\[ \sum_{n=1}^{+\infty} \frac{1}{n} = +\infty \]

\pf{
    \begin{align*}
    & \left( 1+\frac{1}{n} \right)^n<e\\
    & \log\left( 1+\frac{1}{n} \right)^n < \log e \implies \log\left( 1+\frac{1}{n} \right) < \frac{1}{n} \\
    & \underbrace{\log(n+1) - \log n }_{a_{n}}< \frac{1}{n} \\
    & \sum_{k=1}^{n} \log(k+1)-\log k < \sum_{k=1}^{n} \frac{1}{k}\\
    & b_{n} <c_{n}\\
    & \text{Se dimostro che } b_{n} \to +\infty \text{ per il teorema del confronto:} \\
    & \underbrace{\cancel{ \log 2 } - \log 1}_{b_{1}} + \underbrace{{\cancel{ \log 3 } - \cancel{ \log 2 }}}_{b_{2}} + \underbrace{{\cancel{ \log 4 } - \cancel{ \log 3}}}_{b_{3}} + \dots + \log(n+1) - \cancel{ \log n } = b_{n}\\ 
    & b_{n} = \sum_{k=1}^{n} \log(k+1) - \log k = \log(n+1)\\
    & \text{La serie armonica diverge a } +\infty
    \end{align*}
}

\subsubsection{Serie di Mengoli}
\[ \sum_{n=1}^{+\infty} \frac{1}{n(n+1)}=1 \]
\[ \sum_{n=1}^{+\infty} \frac{1}{n}-\frac{1}{n+1} \]
\begin{align*}
    & S_{n}=\sum_{k=1}^{n} \frac{1}{k}-\frac{1}{k+1}=\\
    & = \left( 1 -\cancel{ \frac{1}{2} }\right) + \left( \cancel{ \frac{1}{2} } -\cancel{ \frac{1}{3} } \right) + \left( \cancel{ \frac{1}{3} } -\cancel{ \frac{1}{4} } \right) + \dots \left( \cancel{ \frac{1}{n} } -\frac{1}{n+1} \right)= 1- \frac{1}{n+1}\\
    & \lim_{n \to +\infty} S_n = \lim_{n \to +\infty} \left( 1 - \frac{1}{n+1} \right) = 1
\end{align*}

\subsubsection{Serie Telescopiche}
\[ \sum_{n=1}^{+\infty} (a_{n}-a_{n+1})=a_{1}-\textcolor{orange}{\boxed{\lim_{ n \to +\infty } a_{n+1}}} \]
\[ \textcolor{orange}{\text{La serie e la sua ragione hanno lo stesso carattere: NON limite, ma carattere}} \]
\[ S_{n}=a_{1}-a_{n} \]

\thm{Condizione necessaria per la convergenza}{
    Sia $\displaystyle \sum_{n=1}^{+\infty} a_{n}$ con $a_{n}$ una successione.
    \[ \text{Se } \sum_{n=1}^{+\infty} a_{n} \text{ converge } \implies a_{n}\to 0 \quad (\centernot\impliedby \text{ vedi serie armonica}) \]
}

\pf{
    $a_{n}=S_{n}-S_{n-1}\to S-S = 0$. Non è forma indeterminata perché converge.
}

\thm{Criterio di Convergenza di Cauchy}{
    La serie:
    \[ \sum_{n=1}^{+\infty} a_{n} \text{ conv. } \iff \forall \epsilon> 0 \exists \nu > 0 : |a_{n+1}+a_{n+2}+\dots+a_{n+p}|<\epsilon \quad \forall n > \nu \; \forall p \in \mathbb{N} \]
}

\pf{
    (Per casa)
    $|S_m- S_n| < \epsilon \; \forall n, m > \nu$
    
    \textbf{CONTINUA PER ESERCIZIO}
}

\section{Serie a Termini Non Negativi}

\prop{
    Consideriamo $\displaystyle \left( \sum_{n=1}^{+\infty} a_{n} \; , \; a_{n}\geq 0 \right)$.
    \[ \text{Se }a_{n}\geq 0 \implies \sum_{n=1}^{+\infty} a_{n} \text{ è regolare} \]
}

\pf{
    \[ S_{n}=S_{n-1}+\underbrace{a_{n}}_{\geq 0}\geq S_{n-1} \]
    Allora $S_{n}$ è regolare ($\nearrow$).
}

\ex{
    \begin{itemize}
        \item $\displaystyle \sum_{n=1}^{+\infty} \frac{\sqrt[n]{n!}}{\sqrt{ n }}= +\infty \quad \text{perchè } a_{n} \text{ non tende a 0}$
        \item $\displaystyle \sum_{n=1}^{+\infty} \frac{1}{n}= +\infty \quad \text{anche se } a_{n} \to 0 \text{ (è una condizione necessaria ma non sufficiente)}$
    \end{itemize}
}

\thm{Criterio del confronto (per le serie)}{
    Siano $a_{n},b_{n}\geq 0 \; \forall n$ e sia verificata:
    \[ a_{n}\leq b_{n} \text{ definitivamente} \]
    Allora valgono le seguenti:
    \begin{enumerate}
        \item $\displaystyle \text{Se } \sum_{n=1}^{+\infty} b_{n} \text{ conv.} \implies \sum_{n=1}^{+\infty} a_{n} \text{ conv.}$
        \item $\displaystyle \text{Se } \sum_{n=1}^{+\infty} a_{n} \text{ diverge a } +\infty \implies \sum_{n=1}^{+\infty} b_{n} \text{ diverge a } +\infty$
    \end{enumerate}
}

\pf{
    $S_{n}\leq T_{n} \implies \sum a_{n} \leq \sum b_{n}$
}

\ex{
    \[ \sum_{n=1}^{+\infty} \frac{1}{n^3} \implies \sum_{n=1}^{+\infty} \frac{1}{n^3}<\sum_{n=1}^{+\infty} \frac{1}{n(n+1)} \]
    Dato che $\sum \frac{1}{n(n+1)}$ converge, convergono entrambe.
}

\rmkb{
    \[ \boxed{0<\alpha<1} \]
    \[ \sum_{n=1}^{+\infty} \frac{1}{n^\alpha} > \sum_{n=1}^{+\infty} \frac{1}{n} \]
    Divergono entrambe.
}

\section{Criteri del Confronto Asintotico}

\thm{Criterio del Confronto Asintotico I}{
    Siano $a_{n},b_{n}\geq 0 \; \forall n$.
    \[ \text{Se } \lim_{ n \to \infty } \frac{a_{n}}{b_{n}}=l \in ]0. +\infty[ \]
    Allora le serie $\sum_{n=1}^{+\infty}a_{n}$ e $\sum_{n=1}^{+\infty} b_{n}$ hanno lo stesso carattere.
    Posso scrivere:
    \[ \sum a_{n} \sim \sum b_{n} \]
}

\ex{
    \begin{itemize}
        \item $\displaystyle \sum_{n=1}^{+\infty} \sin \frac{1}{n} = +\infty$
        perchè se $\displaystyle \sum_{n=1}^{+\infty} \frac{1}{n}\implies \lim_{ n \to +\infty } \frac{\sin \frac{1}{n}}{\frac{1}{n}}=1$
        \item $\displaystyle \sum \log\left( 1+ \frac{1}{n} \right) = \sum \sin \frac{1}{n} = +\infty$
        \item $\displaystyle \sum \frac{1}{\sqrt[n]{n!}} = (\text{Come l'armonica?})$
        \item $\displaystyle \sum_{n=1}^{+\infty} \left( 1-\cos \frac{1}{n} \right) \sim \sum_{n=1}^{+\infty} \frac{1}{2n^2} \implies \text{Convergono}$
        \item 
        \begin{align*}
        &\sum_{n=1}^{+\infty} \left( \underbrace{\left( e^{\frac{1}{n}-1} \right)}_{\alpha=1 \; l=1}-\underbrace{\frac{1}{n}}_{{\alpha=1 \; l=1}} \right) \implies \text{Uso Taylor} \implies\sim \sum \frac{1}{n^2} \implies \text{Converge} \implies \\
        & \implies \lim_{ n \to +\infty } \frac{e^{\frac{1}{n}}-1- \frac{1}{n}}{\frac{1}{n^2}} = \lim_{ n \to +\infty } \frac{1+ \frac{1}{n}+ \frac{1}{2n^2} + o\left( \frac{1}{n^2} \right)- 1 - \frac{1}{n}}{\frac{1}{n^2}} =\\
        & =\lim_{ n \to +\infty } \frac{\frac{1}{2n^2}}{\frac{1}{n^2}}=\frac{1}{2} 
        \end{align*}
    \end{itemize}
}

\pf{
    Per Hp:
    \[ \lim_{ n \to +\infty } \frac{a_{n}}{b_{n}} = l > 0 \text{ e } l<+\infty \]
    \[ \forall \epsilon>0 \exists \nu \in \mathbb{N}: l- \epsilon < \frac{a_{n}}{b_{n}} < l + \epsilon \]
    \[ (l-\epsilon)b_{n}<a_{n}<(l+\epsilon)b_{n} \; \forall n > \nu \]
}

\rmkb{
    \[
    \sum_{n=1}^{+\infty} \frac{1}{n^\alpha}=\begin{dcases}
    \alpha\geq 2 \implies \text{Converge} \\
    \text{?} \\
    \alpha=1 \implies \text{Diverge a }+\infty
    \end{dcases}
    \]
}

\prop{
    \textbf{Serie Armonica Generalizzata} \\
    Si consideri la serie:
    \[ \sum_{n=1}^{+\infty} \frac{1}{n^\alpha} \; , \; \alpha \in \mathbb{R} \]
    Allora essa converge $\forall \alpha>1$
    e diverge positivamente $\forall \alpha \leq 1$.
}

\pf{
    (Presente sul Pagani-Salsa) 
    $a_{n}=\frac{1}{n^\alpha}$
    $\alpha>1$
    \begin{align*}
    & S_{2^n -1}=a_{1}+a_{2}+\dots+a_{2^n -1}=\\
    &=1 + \underset{ <2 \cdot \frac{1}{2^\alpha} }{ \textcolor{red}{\boxed{\frac{1}{2^\alpha} +\frac{1}{3^\alpha}}} } +\underset{ < 4 \cdot \frac{1}{4^\alpha} }{ \textcolor{orange}{\boxed{\frac{1}{4^\alpha} +\frac{1}{5^\alpha} +\frac{1}{6^\alpha} +\frac{1}{7^\alpha}}} }+\dots+\\
    & + \frac{1}{(2^{n-1})^\alpha} +\frac{1}{(2^{n-1}+1)^\alpha}+\dots+\frac{1}{(2^{n}-1)^\alpha}
    \end{align*}
    La maggioro $\implies$
    \[ < 1 + \textcolor{red}{\frac{1}{2^{(\alpha-1)}}}+\textcolor{orange}{\frac{1}{2^{2(\alpha-1)}}} +\frac{1}{2^{3(\alpha-1)}}+\dots +\frac{1}{2^{(n-1)(\alpha-1)}}= \sum_{n=0}^{+\infty} \left( \frac{1}{2^{(\alpha-1)}} \right)^n \implies \text{Converge} \]
}

\thm{Criterio Del Confronto Asintotico II}{
    Siano $a_{n}\geq 0, b_{n}>0$. Allora si hanno le seguenti:
    \begin{enumerate}
        \item $\displaystyle \lim_{ n \to +\infty } \frac{a_{n}}{b_{n}}=0$
        (Utile se $a_{n}\to 0$ e $b_{n}\to 0$)
        \begin{itemize}
            \item Se $\sum b_{n}$ converge $\implies \sum a_{n}$ converge
            \item Se $\sum a_{n}$ diverge $\implies \sum b_{n}$ diverge
            \item Esempio b: $\sum \frac{1}{n} =+\infty \implies \sum \frac{1}{\sqrt{ n }}=+\infty$
        \end{itemize}
        \item $\displaystyle \text{Se } \lim_{ n \to +\infty } \frac{a_{n}}{b_{n}} = +\infty$
        Allora: 
        \begin{itemize}
            \item Se $\sum b_{n}$ Div. $\implies \sum a_{n}$ Div.
            \item Se $\sum a_{n}$ Conv. $\implies \sum b_{n}$ Conv.
        \end{itemize}
    \end{enumerate}
}

\pf{
    (Dimostrazione 1)
    \[ \forall \epsilon>0 \exists \nu>0 : 0 \leq \frac{a_{n}}{b_{n}}<\epsilon \quad \forall n > \nu \]
    Scelto $\epsilon=1$ allora $\exists \nu>0:$
    \[ 0\leq a_{n}< b_{n} \quad \forall n>\nu \]
}

\ex{
    \[ \sum_{n=2}^{+\infty} \frac{1}{n^2 \log n} \]
    \[ \sum b_{n}= \sum \frac{1}{n^2} \]
    \[ \lim_{ n \to +\infty } \frac{\cancel{ n^2 }}{\cancel{ n^2 } \log n}=0 \]
}

\subsection{Teorema Criterio degli Infinitesimi}

\thm{Teorema Criterio degli Infinitesimi}{
    Sia $a_{n}\geq 0$ e tale che: 
    \[ \lim_{ n \to +\infty } n^\alpha a_{n} = l\in[0,+\infty] \]
    Allora:
    \begin{enumerate}
        \item Se $0<l<+\infty \implies$
        \begin{itemize}
            \item $\alpha>1\implies \sum a_{n}$ converge.
            \item $\alpha\leq1\implies \sum a_{n}$ diverge.
        \end{itemize}
        \item Se $l=0$ e $\alpha>1 \implies \sum a_{n}$ converge.
        \item Se $l=+\infty$ e $\alpha \leq 1 \implies \sum a_{n}$ diverge.
    \end{enumerate}
}

\ex{
    \begin{itemize}
        \item $\displaystyle \sum_{n=1}^{+\infty} \frac{1}{\sqrt{ n }+2}=+\infty \sim \sum_{n=1}^{+\infty} \frac{1}{\sqrt{ n }} = +\infty \; \text{Perchè } \alpha=\frac{1}{2}$
        \item $\displaystyle \sum_{n=1}^{+\infty} \underbrace{ \sqrt{ n^3 +3n }- \sqrt{ n^3 +2 } }_{a_{n}}$
        \[ \lim_{ n \to +\infty } \frac{n^3 +3n - ( n^3 +2 )}{\sqrt{ n^3 +3n }+ \sqrt{ n^3 +2 }}=\frac{3n-2}{\sqrt{ n^3 +3n } + \sqrt{ n^3 +2 }}= \frac{n(\dots)}{n^{\frac{3}{2}}(\dots)}= \frac{1}{\sqrt{ n }}=0 \]
        \[ \text{Allora } \sum_{n=1}^{+\infty} a_{n} \sim \sum_{n=1}^{+\infty} \frac{1}{\sqrt{ n }}=+\infty \]
    \end{itemize}
}

\subsection{Esempi che non rispettano le Hp precedenti}
\begin{itemize}
    \item $\displaystyle \sum_{n=1}^{+\infty} \frac{n!}{n^n}$
    \item $\displaystyle \sum_{n=1}^{+\infty} \frac{n^n}{3^{n^2}}$
    \item $\displaystyle \sum_{n=1}^{+\infty} \frac{2^n}{n!}$
    \item $\displaystyle \sum_{n=1}^{+\infty} \frac{(2n)!}{(2n)^n}$
\end{itemize}

\subsection{Criterio della Radice per le Serie}

\thm{Criterio della Radice per le Serie}{
    Sia $a_{n}\geq 0$.
    Allora se
    \[ \lim_{ n \to +\infty } \sqrt[n]{a_{n}}=l \]
    Si ha:
    \begin{enumerate}
        \item se $l>1 \implies \sum a_{n}$ diverge a $+\infty$
        \item se $l<1 \implies \sum a_{n}$ converge
    \end{enumerate}
}

\pf{
    \begin{enumerate}
        

    \item Se $l>1$ prendo $\epsilon>0: \quad l-\epsilon>1$.
    Per la definizione di limite:
    \[ \exists \nu>0 : (l- \epsilon) < \sqrt[n]{a_{n}} < (l+\epsilon) \]
    Elevo ad n:
    \[ \exists \nu>0 : (l- \epsilon)^{\textcolor{orange}{n}} < a_{n} < (l+\epsilon)^{\textcolor{orange}{n}} \]
    
    
    \item Sia $0\overset{\text{T.p.s}}{ \leq } l < 1$.
    Sia $\epsilon: l +\epsilon < 1$.
    Poiché:
    \[ \lim_{ n \to +\infty } \sqrt[n]{a_{n}}=l \]
    allora risulta $\exists \nu>0:$
    \[ 0\leq \sqrt[n]{a_{n}}< l +\epsilon\implies 0\leq a_{n} < (l+\epsilon)^n \; \forall n > \nu \]
    Poiché $l+\epsilon<1\implies \sum(l+\epsilon)^n$ converge.
    La dimostrazione segue dal teorema del confronto per le serie.
 
    \end{enumerate}
}

\ex{
    \[ \sum_{n=1}^{+\infty} \frac{n!}{n^n} \quad a_{n}\geq 0 \]
    Criterio della Radice:
    \[ \lim_{ n \to +\infty } \sqrt[n]{\frac{n!}{n^n}}=\lim_{ n \to +\infty } \frac{\sqrt[n]{n!}}{n}=\frac{1}{e} < 1 \]
}

\thm{Criterio del Rapporto per le Serie}{
    Sia $a_{n}> 0$ e tale che esista:
    \[ \lim_{ n \to +\infty } \frac{a_{n+1}}{a_{n}}=l \]
    Allora:
    \begin{enumerate}
        \item $\text{Se } l>1\implies \sum_{n=1}^{+\infty} a_{n}=+\infty$
        \item $\text{Se }l<1 \implies \sum_{n=1}^{+\infty} a_{n} \text{ converge}$
    \end{enumerate}
}

\pf{
    Segue dal Criterio Rapporto radice per le Successioni.
}

\subsubsection{Costante di Eulero Mascheroni}

\thm{Costante di Eulero Mascheroni}{
    La successione:
    \[ \left( \sum_{k=1}^{n} \frac{1}{k} \right)-\log n \]
    converge ad un numero finito non nullo che chiameremo:
    \[ \gamma = \text{ costante di Eulero Mascheroni} \approx 0,577 \]
    \[ \lim_{ n \to +\infty } \frac{S_{n}}{\log n}=\lim_{ n \to +\infty } \frac{S_{n}-\log n +\log n}{\log n}=\lim_{ n \to +\infty } \frac{S_{n}-\log n}{\log n}+1=1 \]
}

\section{Serie a Segni Alterni}

\defn{Serie a Segni Alterni}{
    Sia $a_{n}\geq 0\; \forall n$, la serie:
    \[ \sum_{n=1}^{+\infty} (-1)^{n-1} a_{n}=a_{1}-a_{2}+a_{3}-a_{4}+\dots \]
    Si chiama \textbf{serie a segni alterni} di termine generale $a_{n}$.
}

\ex{
    \begin{itemize}
        \item $\displaystyle \sum_{n=1}^{+\infty} \frac{(-1)^{n-1}}{n}=1 -\frac{1}{2}+\frac{1}{3}-\frac{1}{4}+\dots \qquad \text{Serie armonica a segni alterni}$
        \item $\displaystyle \sum_{n=1}^{+\infty} (-1)^n \arcsin\left( \frac{1}{n} \right)= -\arcsin 1 +\arcsin \frac{1}{2} - \dots$
    \end{itemize}
}

\subsection{Criterio di Leibniz}

\thm{Criterio di Leibniz}{
    \begin{enumerate}
        \item Sia $a_{n}\geq 0$
        \item Sia $a_{n} \searrow$ (decrescente)
        \item Sia $a_{n}$ infinitesima $\left(\displaystyle \lim_{ n \to +\infty }a_{n}=0\right)$
    \end{enumerate}
    \[ \text{Allora } \sum_{n=1}^{+\infty} (-1)^{n-1}a_{n} \text{ converge} \]
    Inoltre, detta $S$ la sua somma, si ha la seguente:
    \[ |S-S_{n}|\leq a_{n+1} \quad (\text{Stima del resto}) \]
}

\pf{
    Considero $S_{2n}$ e $S_{2n+1}$.
    Dimostro che sono convergenti allo stesso limite.
    1) Dimostriamo che $S_{2n}$ e $S_{2n+1}$ sono monotone e limitate.
    \[ S_{\underbracket{2n+2}}=S_{2n} \underbrace{+ a_{2n+1}-a_{2n+2} }_{{\geq 0 \; (a_{n} \searrow)}}\geq S_{\underbracket{2n}} \quad \nearrow \]
    \[ S_{\underbracket{2n+1}}=S_{2n-1}\underbrace{ - a_{2n}+a_{2n+1} }_{{\leq 0}}\leq S_{\underbracket{2n-1}} \quad \searrow \]
    $S_{2n} \nearrow \quad S_{2n+1}\searrow$
    Allora:
    \[ S_{3}\geq S_{2n+1}\geq S_{2n}\geq S_{2}\implies\text{ quindi } S_{2n} \text{ e } S_{2n+1} \text{ sono anche limitate.} \]
    $\implies$ sono convergenti per il Teorema delle Successioni monotone.
    Ma poiché:
    \[ S_{2n+1}-S_{2n}=a_{2n+1} \]
    Al limite per $n\to + \infty$:
    \[ S_{d}-S_{p}= 0 \implies S_{d}=S_{p}=\boxed{S} \]
    \[ \inf \{S_{2n+1}\} = S = \sup \{S_{2n}\} \]
    \[
    |S-S_{n}|=\begin{dcases}
    n \text{ pari} \\
    n \text{ dispari}
    \end{dcases}
    \]
    \begin{align*}
    &|S-S_{2n}| \underset{\substack{\text{(perché } S \\ \text{è il sup)}}}{=} S- S_{2n} \\
    & S - S_{2n} \underset{S \text{ è inf}}{\leq} S_{2n+1}-S_{2n}=a_{2n+1}\\
    & |S-S_{2n+1}| \underset{S \text{ è inf}}{=} S_{2n+1}-S \leq S_{2n+1}-S_{2n+2}=+a_{2n+2}
    \end{align*}
}

\ex{
    \[ \sum_{n=1}^{+\infty} \frac{(-1)^{n-1}}{n} \; \text{ Converge} \]
}

\subsection{Assoluta Convergenza}

\defn{Successione Assolutamente Convergente}{
    Sia $a_{n}$ una successione.
    La serie $\displaystyle \sum_{n=1}^{+\infty} a_{n}$ si dice \textbf{assolutamente convergente} se converge:
    \[ \sum_{n=1}^{+\infty} |a_{n}| \]
}

\rmkb{
    La convergenza finora citata è detta \textbf{convergenza semplice}.
}

\prop{
    Se $\displaystyle \sum_{n=1}^{+\infty} a_{n}$ converge assolutamente (ovvero se $\sum |a_{n}|$ converge)
    \[ \implies \sum_{n=1}^{+\infty} a_{n} \text{ converge} \]
    ($\centernot\impliedby$ Esempio: Armonica a segni alterni)
}

\ex{
    \[ \sum_{n=1}^{+\infty} \left|\frac{(-1)^n - n +1}{n^3+2n}\right|\leq \sum \frac{n+2}{n^3 +2n} \sim \sum \frac{1}{n^2} \]
}

\pf{
    \[ \Big||a_{n+1}|+|a_{n+2}|+\dots+|a_{n+p}|\Big|<\epsilon \quad \forall n>\nu \quad \forall p \in \mathbb{N} \]
    \[ \Big||a_{n+1}+a_{n+2}+\dots+a_{n+p}|\Big|<\Big||a_{n+1}|+|a_{n+2}|+\dots+|a_{n+p}|\Big| \]
    \[ \implies\left|a_{n+1}+a_{n+2}+\dots+a_{n+p}\right|<\Big||a_{n+1}|+|a_{n+2}|+\dots+|a_{n+p}|\Big| \]
}